\section{One pointwise completion theorem}
\label{simp:completion-section}

The completion argument uses only one potential.  Its charge coefficient
is one, exactly matching the unit of spare column capacity in a cost-two
NEW label.  No intermediate capacity surcharge is needed.

For a capped factor with parameter $\theta$, and for an attenuated factor
with parameters $(a,\theta)$, respectively, define
\[
 F_{C,\theta}(r)=\max(1,\theta r),\qquad
 F_{a,\theta}(r)=\max\{1,\theta\max_{s\in\mathbb Z_{\ge0}}a^s(r+s)\}.
\]
The fixed parameters obey $\theta a^s(s+1)\le1$ for every integer $s\ge0$.
Comparing consecutive terms reduces this condition to the single index
$s=\lfloor a/(1-a)\rfloor$.  The maximum defining $F_{a,\theta}$ is
finite: the ratio of consecutive nonconstant terms is
$a(r+s+1)/(r+s)$, eventually below one.  In fact indices above
$\lceil a/(1-a)\rceil$ can be discarded uniformly for $r\ge0$.

Every $F_c$ is nondecreasing, $F_c(r)/r$ is nonincreasing, and $F_c(r)=1$
for $0\le r\le1$.  Moreover
\[
 aF_{a,\theta}(r+1)\le F_{a,\theta}(r).
\]
Set
\[
 \mathcal F(r)=\prod_c F_c(r)^{\omega_c},\qquad
 \chi(\delta)=\log\mathcal F(e^\delta),\qquad
 \Gamma=\prod_c\theta_c^{\omega_c}.
\]
Then $\chi$ is convex, nondecreasing, one-Lipschitz, zero on
$(-\infty,0]$, and $0\le\chi(\delta)\le\delta^+$.
Indeed each $\log F_c(e^\delta)$ is a maximum of zero and functions
$\log(\theta a^s)+\log(e^\delta+s)$, all convex; its slopes lie in $[0,1]$.
Also $\mathcal F(r)/r\to\Gamma$ as $r\to\infty$.

\begin{lemma}[Unit charge]\label{simp:unit-charge}
For every real $\delta$,
\[
 \min\{\log2,\chi(\delta)\}\le\delta^+-\chi(\delta).
\]
\end{lemma}
\begin{proof}
Three uses of the attenuation recursion give $F_{a,\theta}(4)\le a^{-3}$.
For the three capped factors, the respective values at four are $4,1,2$.
The exact weights in the score table satisfy the two rational inequalities
\[
 w:=2\omega_{C,1}+\omega_{C,1/2}<\frac23,
 \qquad s:=3\sum_{c\text{ attenuated}}\omega_c\frac{1-a_c}{a_c}<\frac16.
\]
Since $-\log a\le(1-a)/a$ and $\log2\ge1/2$,
\[
 \chi(\log4)\le w\log2+s<\tfrac23\log2+\tfrac16\le\log2.
\]
Convexity and $\chi(0)=0$ imply $\chi(\delta)\le\delta/2$ on
$[0,\log4]$.  Above $\log4$, one-Lipschitzness implies
$\delta-\chi(\delta)\ge\log4-\chi(\log4)>\log2$.
For $\delta\le0$ both sides vanish.
\end{proof}

For $p\ge0$, $\alpha>0$, $\max(X,p)\le V\le X+p$, put $b=\max(p,\alpha)$.
The ordinary kernel used in the proof is
\begin{equation}\label{simp:strong-kernel}
 \widehat K_{p,\alpha}(X,V)=
 \begin{cases}
 V\mathcal F(\alpha/V),&V\le b,\\
 V\mathcal F(\alpha/X)/\mathcal F(\alpha V/(bX)),&V\ge b.
 \end{cases}
\end{equation}
The formulas agree at $V=b$; continuous limits specify zero arguments.
This kernel is used only in the analysis of the final matching, not as an
extra option in the auxiliary algorithm.  The private and NEW kernels are
\[
 B_\alpha(X)=\prod_c\max(X,\theta_c\alpha)^{\omega_c},
 \qquad W_p(V)=\min(V,2p).
\]

\begin{lemma}[Ordinary envelope]\label{simp:ordinary-envelope}
The original physical kernel $\mathcal H_{p,\alpha}(X,V)$ is at most
$\widehat K_{p,\alpha}(X,V)$, and
\[
 \widehat K_{p,\alpha}(X,V)=
 \inf_{\substack{u>0,\ u\ge X,\ w\ge1\\uw\ge V}}
 uw\,\frac{\mathcal F(\alpha/u)}
 {\mathcal F(w\min(1,\alpha/p))}.
\]
At $p=0$, interpret $\min(1,\alpha/p)$ as one.
\end{lemma}
\begin{proof}
To prove the identity put $v=uw$ and $r=\alpha/u$.  The minimized expression
is $v\mathcal F(r)/\mathcal F((v/b)r)$ with $X\le u\le v$.
For $v\ge b$, convexity of $\chi$ makes the ratio nonincreasing in $r$,
so $u=X$ minimizes it.  For $v\le b$, the ratio is nondecreasing in $r$,
so use $u=v$; then $\mathcal F(\alpha/b)=1$.
In both cases the resulting expression is nondecreasing in $v$, by
monotonicity of $\mathcal F(r)/r$.  Its minimum is therefore at $v=V$.

For domination fix a feasible $(u,w)$ and let
$r=\alpha/u$, $q=w\min(1,\alpha/p)$.  Thus $w\ge\max(1,q)$ and
\[
 p/u\le w\min(1,r/q),\qquad V/u\le\min(w,1+rw/q),\qquad X/u\le1.
\]
For each factor, monotonicity of $F$ and reverse monotonicity of $F(r)/r$
bound its singleton, raw, and private-floor entries by $wF(r)/F(q)$.
For a capped entry, $\min(V,X/\theta)/u\le\min(w,1/\theta)
\le w/F(q)$.  For an attenuated entry it remains to show
\[
 a\min(1,w^{-1}+r/q)F(q)\le F(r).
\]
If $q\le1$, use $F(q)=1$.  If $1\le q\le r+1$, use
$aF(q)\le aF(r+1)\le F(r)$.  If $q>r+1$, use
$a(r+1)F(q)/q\le aF(r+1)\le F(r)$.
Take the weighted geometric mean and then the infimum.
\end{proof}

\begin{theorem}[Common cost-two completion]\label{simp:completion}
Fix disjoint residual bundles and reservation data
$p_{ij}\le\alpha_i e^{\rho_j}$, $p_{ii}=\alpha_i e^{\rho_i}>0$, $\rho_j\ge0$.
Let $Q$ be the maximum utility product after a bijective rematching of the
reserved items.  Let $a_{ij},z_{ij},b_i,d_i\ge0$ obey
\[
 d_i+b_i+\sum_j(a_{ij}+z_{ij})=1,
 \qquad b_j+\sum_i(a_{ij}+2z_{ij})\le1.
\]
With $X_i=f_i(S_i)$ and $V_{ij}=f_i(S_i+j)$,
\[
 Q\ge
 \prod_{ij}\widehat K_{ij}^{a_{ij}}
 \prod_{ij}W_{ij}^{z_{ij}}
 \prod_iV_{ii}^{b_i}\prod_iB_{\alpha_i}(X_i)^{d_i}.
\]
Zero exponent factors are omitted; zero supported factors make the assertion
immediate.  The theorem includes any mixture of OLD, NEW, full HOME, and
private factors on the same residual outcome.
\end{theorem}
\begin{proof}
Normalize an optimal dual of the actual final perfect matching as
$u_i>0$, $\pi_j\ge0$, $\min_j\pi_j=0$, so $u_i\ge X_i$ and
$Q=\prod_i u_i\prod_j e^{\pi_j}$.
Write $r_i=\alpha_i/u_i$, $\delta_i=\log r_i$, and $\chi_i=\chi(\delta_i)$.
The HOME and reservation inequalities imply
$r_j\le e^{\pi_j-\rho_j}\le e^{\pi_j}\min(1,\alpha_i/p_{ij})$.
The envelope lemma therefore gives
\[
 \log(\widehat K_{ij}/u_i)\le\pi_j+\chi_i-\chi_j,
 \qquad \log(B_{\alpha_i}(X_i)/u_i)\le\chi_i.
\]
The full HOME value satisfies $\log(V_{ii}/u_i)\le\pi_i$ directly.
The singleton and union bounds give
\[
 \log(W_{ij}/u_i)\le
 \pi_j-(\delta_j-\delta_i-\log2)^+.
\]
Let $c_j=b_j+\sum_i(a_{ij}+z_{ij})$.  Subtracting $\log Q$ from the proposed
log product yields an upper bound
\[
 -\sum_j(1-c_j)(\pi_j-\chi_j)
 +\sum_{ij}z_{ij}
 [\chi_j-\chi_i-(\delta_j-\delta_i-\log2)^+].
\]
Every bracket is at most $\min(\log2,\chi_j)$, using nonnegativity and
one-Lipschitzness of $\chi$.  The unit-charge lemma and $r_j\le e^{\pi_j}$
bound it by $\delta_j^+-\chi_j\le\pi_j-\chi_j$.
Finally $1-c_j\ge\sum_i z_{ij}$ by the cost-two column constraint, so the
entire expression is nonpositive.
\end{proof}
